\documentclass[12pt, letterpaper, onecolumn]{book}

% --- Preamble: Packages and Settings ---
\usepackage[utf8]{inputenc} % Set input encoding to UTF-8
\usepackage[T1]{fontenc}    % Use T1 font encoding
\usepackage{amsmath}        % For advanced math environments
\usepackage{graphicx}       % For including images
\usepackage[margin=1in]{geometry} % Set margins to 1 inch
\usepackage{enumitem}


% Optional: For better spacing and line breaks in the bibliography
\usepackage{url}
\usepackage{hyperref} % For clickable links in the PDF

% Set document title, author, and date
\title{Large Etendue Synoptic Search Systems and Models}
\author{Earl Spillar}
\date{\today}

% --- Document Start ---
\begin{document}

% Generates the title, author, and date
\maketitle

\tableofcontents

\chapter{Introduction}
Capturing the thinking and mathematics behind using Large \'etendue optical
systems for synoptic search.  
Although evolving, this is meant to be more of a reference work, and perhaps a treatise, 
as opposed to be working document: although it will probably start that way!

\chapter{Equations and Relations}
% --- Main Body ---
\section{Search Rate}
\subsection{Variables}
\begin{enumerate}[noitemsep]
	\item $\Omega$ Field of view 
	\item $\omega$ Effective fov of the system, 
	something like the maximum of the blur size and the pixel size.
	\item $\alpha$ brightness of the object in units like
	photons per square meter per second
	\item $\beta$ brightness of the object in units like
	photons per square meter per steradian
	\item $\gamma$ desired signal to noise ratio
	\item $\eta$ efficiency, including a large number of other factors (see below))
	\item T effective net dwell time.
	\item t integration time, $t \le T.$
%	\item $\dot{\theta}$
\end{enumerate}

%Factors included in $\eta$ include:
%\begin{itemize}[noitemsep]
%	\item quantum efficiency,
%	\item ``shutter open" time,
%	\item atmospheric transmission,
%	\item light outside of the photometry aperture into other %pixels,
%\end{itemize}


\subsection{Derivation}
\subsection{Search rate}
Under the assumption that we have included $\eta$ properly
and that we can neglect read noise and dark current (background limited 
performance or BLIP) the search rate
$$ SR = {1 \over T} = \frac{\Omega A \eta \alpha^2 }{\beta \gamma^2 } $$
Which is sometimes given the unit of Hershels, square meters square degrees
(I think this is wrong, need to check it!)

\subsection{Maximum Integration Time]


\subsection{Background}
Letter paper is a standard size for documents in the U.S. and Canada, measuring $8.5 \times 11$ inches. By using the $\backslash$usepackage[margin=1in]{geometry} package, we ensure the text fits nicely on the page.

\section{Methods}
This section describes the methodology. Here is an example of an equation:
$$
E = mc^2
$$
You can use inline math like $\sqrt{x^2+y^2}$ within your text.

\section{Results and Discussion}
To include an image, place the file (e.g., `figure.png`) in the same folder as your `.tex` file and use the `figure` environment.

%\begin{figure}[h]
%    \centering
%    \includegraphics[width=0.8\textwidth]{figure.png}
%    \caption{A descriptive caption for your figure.}
%    \label{fig:myfig}
%\end{figure}

For references, you can use the $\backslash$cite command and a bibliography style. For instance, see a general reference \cite{Knuth:1984}.

\section*{Conclusion}
OK

% --- Bibliography ---

\chapter{Appendices}
\section{}

\begin{thebibliography}{9}

\bibitem{Knuth:1984}
Donald E. Knuth.
\emph{The \TeX{}book}.
Addison-Wesley Professional, 1984.

\end{thebibliography}

\end{document}
